% Gemini theme
% https://github.com/anishathalye/gemini

\documentclass[final]{beamer}

% ====================
% Packages
% ====================

\usepackage[T1]{fontenc}
\usepackage{lmodern}
\usepackage[size=custom,width=60.96,height=91.44,scale=1.0]{beamerposter}
\usetheme{gemini}
\usecolortheme{harvard}
\usepackage{graphicx}
\usepackage{booktabs}
\usepackage{tikz}
\usepackage{pgfplots}
\usepackage{caption}
\pgfplotsset{compat=1.14}
\usepackage{anyfontsize}

% ====================
% Lengths
% ====================

% Poster: 24in x 36in (60.96cm x 91.44cm)
% 2 columns → choose widths such that:
% (2+1)*sepwidth + 2*colwidth = \paperwidth

\newlength{\sepwidth}
\newlength{\colwidth}

\setlength{\sepwidth}{0.02\paperwidth}
\setlength{\colwidth}{0.47\paperwidth}


\newcommand{\separatorcolumn}{\begin{column}{\sepwidth}\end{column}}

% ====================
% Title
% ====================

\title{Improving GCG}
\logoright{\includegraphics[width=6cm]{qrcode.png}}

\author{Ege Çakar \and Hannah Guan \and Kayden Kehe}

\institute[shortinst]{Harvard School of Engineering and Applied Science}

% ====================
% Footer (optional)
% ====================

\footercontent{
  \href{github.com/Ege-Cakar/ImprovingGCG}{github.com/Ege-Cakar/ImprovingGCG} \hfill
  Harvard CS2881 - AI Safety and Alignment}
% (can be left out to remove footer)

% ====================
% Logo (optional)
% ====================

% use this to include logos on the left and/or right side of the header:
% \logoright{\includegraphics[height=7cm]{logo1.pdf}}
% \logoleft{\includegraphics[height=7cm]{logo2.pdf}}

% ====================
% Body
% ====================

\begin{document}

% Refer to https://github.com/k4rtik/uchicago-poster
% logo: https://www.cam.ac.uk/brand-resources/about-the-logo/logo-downloads
\addtobeamertemplate{headline}{}
{
    \begin{tikzpicture}[remember picture,overlay]
      \node [anchor=north west, inner sep=3cm] at ([xshift=0.0cm,yshift=1.5cm]current page.north west)
      {\includegraphics[height=6cm]{logo/logo.png}};
      
    \end{tikzpicture}
}


\begin{frame}[t]
\begin{columns}[t]
\separatorcolumn

\begin{column}{\colwidth}

  \begin{alertblock}{Problem \& Motivation}
    \begin{itemize}
        \item Modern LLMs show brittle alignment despite strong safety training.
        \item Greedy Coordinate Gradient (GCG)-based search has been shown to effectively jailbreak aligned models.
        \item But GCG is slow and targets the output instead of the internal refusal mechanism.
    \end{itemize}
    We ask: \textit{Can we improve GCG by explicitly targeting a model's internal refusal direction and making its optimization more efficient?}
  \end{alertblock}

  \begin{block}{Activation-Guided GCG}
    
    Goal: optimize an adversarial suffix using GCG to \textbf{target the model's internal refusal direction}.

    \textbf{Refusal direction:} Per-layer refusal directions (Arditi et al.\ 2024) are calculated as
    $$r^{(\ell)} = \mu^{(\ell)}_{\textrm{harmful}} - \mu^{(\ell)}_{\textrm{harmless}}$$
    where $\mu^{(\ell)}$ is the mean activation at layer $\ell$. We use the precomputed direction at layer $\ell^*=14$, position $p=-1$ --- the (layer, position) pair Arditi et al.\ identify as most reducing refusal under ablation.

    \textbf{Activation-based losses:} We develop a family of activation-based losses which encourage the suffix to suppress the model's refusal direction:
    \begin{itemize}
        \item[-] Single: Minimize (squared) projection at the most influential (layer, position) pair.
        \item[-] Negative: Push the projection at the most influential (layer, position) pair to be negative (actively oppose refusal).
        \item[-] Layer: Minimize projections across all token positions at the most influential layer.
        \item[-] Token: Minimize projections at a single token position across all layers.
        \item[-] All: Minimize projections across all layers and all positions.
    \end{itemize}


    \vspace{0.5em}
    

  \end{block}


  \begin{alertblock}{Soft GCG}
      Goal: accelerate GCG by relaxing the discrete token search into a \textbf{continuous optimization over vocabulary logits}, closing the resulting ``projection gap'' via temperature annealing.

      \vspace{0.3em}

      \textbf{Mechanism.} For each of the $L$ suffix positions, maintain a learnable logit vector $\phi_i \in \mathbb{R}^{|V|}$ over the vocabulary. At every step, form a \emph{soft} one-hot via the Gumbel--Softmax (Jang et al.\ 2017):
      $$\tilde{s}_i = \text{Softmax}\!\left(\frac{\phi_i + g_i}{\tau}\right), \qquad g_i \sim \text{Gumbel}(0,1),$$
      and embed it as a probability-weighted mix of vocabulary embeddings, $e(\tilde{s}_i) = \tilde{s}_i^{\top} E$. This is fully differentiable, so the adversarial loss backpropagates through to the logits $\phi$. After $T$ Adam steps on $\phi$, discretize by $s_i = \arg\max_v \phi_i[v]$; the resulting hard suffix can be used directly or passed as initialization to discrete GCG.

      \vspace{0.3em}

      \textbf{Temperature annealing.} $\tau$ starts high to allow broad exploration across tokens and cools toward zero so $\tilde{s}_i$ collapses onto a single token; we compare a linear schedule ($\tau: 2.0 \to 0.1$) against a 3-phase ``slushy'' schedule (exploration $2.5\!\to\!1.0$, refinement $1.0\!\to\!0.5$, solidification $0.5\!\to\!0.01$).
    \begin{figure}
        \centering
        \includegraphics[width=0.4\linewidth]{soft_opt_schedule.png}
    \end{figure}

  \end{alertblock}

  \begin{block}{Evaluation Metrics}

    \textbf{How do we define ``harmfulness" for evaluation?}
    \begin{itemize}
        \item Substring ASR: Labels a completion as harmful if it does \textit{not} contain any of a fixed list of refusal-style prefixes (e.g., ``I'm sorry", ``I cannot").
        \item LlamaGuard2 ASR: Labels a completion as harmful if it is labeled unsafe by the LlamaGuard2 safety classifier.
        \item HarmBench ASR: Labels a completion as harmful if it is labeled unsafe by the HarmBench evaluation framework.
    \end{itemize}

  \end{block}

\end{column}

\separatorcolumn

\begin{column}{\colwidth}

  \begin{alertblock}{Experiment Settings}
      \textbf{Datasets:} Activation-Guided GCG trains on AdvBench goals and evaluates on the harmful-test split from Arditi et al.\ (2024). Soft GCG evaluates on AdvBench (paired goal/target strings); sweep runs use 25 train + 10 test prompts.

      \vspace{0.5em}

      \textbf{Activation-Guided GCG:} Target \texttt{Llama-2-7b-chat-hf}; suffix length 20; 200 discrete optimization steps; standard \texttt{llm\_attacks} parameters; refusal direction at $\ell^*=14$, $p=-1$.

      \vspace{0.5em}

      \textbf{Soft GCG:} Targets \texttt{Llama-2-7b-chat-hf} and the Gemma 3 family (270M, 1B, 4B, 12B); suffix length 20; Adam LR $0.1$ on suffix logits. Llama 2 sweep: Pure GCG (500), Pure Soft (500), and hybrid (e.g., 250/250, 400/100), with 128 GCG candidates per step; long-run seed comparison uses 2000 iterations. Gemma: 500 steps of Pure Soft, batch size 10. Objectives: cross-entropy, or a Carlini--Wagner margin loss ($\kappa=5.0$, first-token weight $w_0=5.0$) for more robust discretized suffixes.
  \end{alertblock}


  \begin{block}{Experiment Results}

        \vspace{0.25}
  
      \textbf{Activation-based objectives outperform standard GCG at 200 steps.} The ``All'' objective (global suppression across all layers and positions) reaches \textbf{0.91 substring ASR vs.\ 0.76 for GCG}, approaching the direction-ablation ceiling of 0.98. Single, Layer, and Negative all match or exceed GCG; gains carry over to LlamaGuard2 and HarmBench ASR (max 0.09 and 0.04 vs.\ 0.00 for GCG).
      \begin{figure}[h]
        \centering
        \includegraphics[width=0.95\columnwidth]{activation_asr.png}
      \end{figure}

    \vspace{0.25em}

    \textbf{Soft GCG matches GCG's performance on Llama-2-7B-Chat with a 33$\times$ wall-clock speedup} (2.5 min vs.\ 81 min on the same hardware; up to 43$\times$ in shorter 500-step runs). No discrete-GCG polish is required.
    \begin{figure}[h]
        \centering
        \includegraphics[width=0.95\columnwidth]{sgcg_llama2.png}
    \end{figure}

    \vspace{0.25em}

    \textbf{Soft GCG attacks the Gemma 3 family with success that decays sharply with model size:} 100\% on Gemma3-270M$^*$, 57.7\% on Gemma3-1B, 33.6\% on Gemma3-4B, and 0\% on Gemma3-12B (substring ASR; $^*$270M became incoherent after optimization). Larger and better safety-trained models resist both activation- and suffix-based attacks at our compute budget.
    \begin{figure}[h]
        \centering
        \includegraphics[width=0.95\columnwidth]{sgcg_gemma.png}
    \end{figure}

  \end{block}

\end{column}

\separatorcolumn
\end{columns}
\end{frame}

\end{document}

